\begin{thebibliography}{99}

\bibitem{ArratiaGoldsteinGordon1989}
R. Arratia, L. Goldstein and L. Gordon,
\emph{Two moments suffice for Poisson approximations: the Chen--Stein method},
Ann. Probab. \textbf{17} (1989), no.~1, 9--25.
\href{https://doi.org/10.1214/aop/1176991491}{doi:10.1214/aop/1176991491}.

\bibitem{ArratiaGoldsteinGordon1990}
R. Arratia, L. Goldstein and L. Gordon,
\emph{Poisson approximation and the Chen--Stein method},
Statist. Sci. \textbf{5} (1990), no.~4, 403--424.
\href{https://doi.org/10.1214/ss/1177012015}{doi:10.1214/ss/1177012015}.

\bibitem{BalasubramanianShorey1993}
R. Balasubramanian and T. N. Shorey,
\emph{Squares in products from a block of consecutive integers},
Acta Arith. \textbf{65} (1993), no.~3, 213--220.
\href{https://doi.org/10.4064/aa-65-3-213-220}{doi:10.4064/aa-65-3-213-220}.

\bibitem{Barbour1988SteinProcess}
A. D. Barbour,
\emph{Stein's method and Poisson process convergence},
J. Appl. Probab. \textbf{25A} (1988), 175--184.
\href{https://doi.org/10.2307/3214155}{doi:10.2307/3214155}.

\bibitem{Berry1941}
A. C. Berry,
\emph{The accuracy of the Gaussian approximation to the sum of independent
variates},
Trans. Amer. Math. Soc. \textbf{49} (1941), no.~1, 122--136.
\href{https://doi.org/10.1090/S0002-9947-1941-0003498-3}
{doi:10.1090/S0002-9947-1941-0003498-3}.

\bibitem{CheeKiahPurkayasthaWang2013}
Y. M. Chee, H. M. Kiah, P. Purkayastha and C. Wang,
\emph{Cross-bifix-free codes within a constant factor of optimality},
IEEE Trans. Inform. Theory \textbf{59} (2013), no.~7, 4668--4674.
\href{https://doi.org/10.1109/TIT.2013.2252952}{doi:10.1109/TIT.2013.2252952}.

\bibitem{HalterKoch2013}
F.~Halter-Koch,
\emph{Quadratic Irrationals: An Introduction to Classical Number Theory},
Monographs and Textbooks in Pure and Applied Mathematics, vol.~306,
CRC Press, Boca Raton, 2013.
\href{https://doi.org/10.1201/b14968}{doi:10.1201/b14968}.

\bibitem{Krokowski2017}
K. Krokowski,
\emph{Poisson approximation of Rademacher functionals by the Chen--Stein
method and Malliavin calculus},
Commun. Stoch. Anal. \textbf{11} (2017), no.~2, 195--222.
\href{https://doi.org/10.31390/cosa.11.2.05}{doi:10.31390/cosa.11.2.05};
\href{https://arxiv.org/abs/1505.01417v3}{arXiv:1505.01417v3}.

\bibitem{LaishramShorey2004}
S. Laishram and T. N. Shorey,
\emph{Number of prime divisors in a product of consecutive integers},
Acta Arith. \textbf{113} (2004), no.~4, 327--341.
\href{https://doi.org/10.4064/aa113-4-3}{doi:10.4064/aa113-4-3}.

\bibitem{NicolasRobin1983}
J.-L. Nicolas and G.~Robin,
\emph{Majorations explicites pour le nombre de diviseurs de $N$},
Canad. Math. Bull. \textbf{26} (1983), no.~4, 485--492.
\href{https://doi.org/10.4153/CMB-1983-078-5}{doi:10.4153/CMB-1983-078-5}.

\bibitem{PoulyLatticePoissonFlows2026}
B. Pouly,
\emph{Records, scale flows, and Gaussian limits from lattice Poisson fields},
Preprint, 2026, with technical companion.
Concept DOI \href{https://doi.org/10.5281/zenodo.22286837}{\nolinkurl{10.5281/zenodo.22286837}}.

\bibitem{Rollin2012SteinFactors}
A. R\"ollin,
\emph{On the optimality of Stein factors},
in \emph{Probability Approximations and Beyond}, Lecture Notes in Statistics
\textbf{205}, Springer, 2012, 61--72.
\href{https://doi.org/10.1007/978-1-4614-1966-2_5}{doi:10.1007/978-1-4614-1966-2\_5};
\href{https://arxiv.org/abs/0706.0879v3}{arXiv:0706.0879v3}.

\bibitem{Shorey1994BakerApplications}
T. N. Shorey,
\emph{Applications of Baker's theory of linear forms in logarithms to
exponential diophantine equations},
RIMS K\^oky\^uroku \textbf{886} (1994), 48--60.
\href{https://www.kurims.kyoto-u.ac.jp/~kyodo/kokyuroku/contents/pdf/0886-05.pdf}{RIMS archive}.

\bibitem{Tenenbaum2015}
G. Tenenbaum,
\emph{Introduction to Analytic and Probabilistic Number Theory},
3rd ed., Graduate Studies in Mathematics 163, American Mathematical Society, 2015.
\href{https://doi.org/10.1090/gsm/163}{doi:10.1090/gsm/163}.

\end{thebibliography}
